% This LaTeX document needs to be compiled with XeLaTeX.
\documentclass[10pt]{article}
\usepackage[utf8]{inputenc}
\usepackage{ucharclasses}
\usepackage{amsmath}
\usepackage{amsfonts}
\usepackage{amssymb}
\usepackage[version=4]{mhchem}
\usepackage{stmaryrd}
\usepackage{bbold}
\usepackage{polyglossia}
\usepackage{fontspec}
\setmainlanguage{english}
\setotherlanguages{hindi}
\IfFontExistsTF{Noto Serif Devanagari}
{\newfontfamily\hindifont{Noto Serif Devanagari}}
{\IfFontExistsTF{Kohinoor Devanagari}
  {\newfontfamily\hindifont{Kohinoor Devanagari}}
  {\IfFontExistsTF{Devanagari MT}
    {\newfontfamily\hindifont{Devanagari MT}}
    {\IfFontExistsTF{Lohit Devanagari}
      {\newfontfamily\hindifont{Lohit Devanagari}}
      {\IfFontExistsTF{FreeSerif}
        {\newfontfamily\hindifont{FreeSerif}}
        {\newfontfamily\hindifont{Arial Unicode MS}}
}}}}
\IfFontExistsTF{CMU Serif}
{\newfontfamily\lgcfont{CMU Serif}}
{\IfFontExistsTF{DejaVu Sans}
  {\newfontfamily\lgcfont{DejaVu Sans}}
  {\newfontfamily\lgcfont{Georgia}}
}
\setDefaultTransitions{\lgcfont}{}
\setTransitionsForDevanagari{\hindifont}{\rmfamily}

\begin{document}
$L 15$\\
/

L15: Primitive Roots

\section*{The order of an integer.}
\begin{itemize}
  \item Let $m$ be a positive integer and a be an integer relatively prime to $m$. Euler's Theorem tells us that $a^{\phi(m)} \equiv 1(\bmod m)$. However oit may happen that some smaller positive power than $\phi(m)$, say $n$, will result in $a^{n} \equiv 1(\bmod m)$.\\
Definition: Let $a, m \in \mathbb{Z}$ with $m>0$ and $(a, m)=1$. The order of a modulo $m$, denoted ord ${ }_{m}$ a, is the least positive integer $n$ for which $a^{n} \equiv 1(\bmod m)$.\\
Example:- Consider ord, 2 . $2^{1} \equiv 2(\bmod 7), 2^{2} \equiv 4(\bmod 7), 2^{3} \equiv 1(\bmod 7)$.\\
Thus ord, $2=3$.\\
Euler's Theorem guarantees that
\end{itemize}

$$
\text { ord, } 2 \leqslant \phi(2)=6 \text {. }
$$

\begin{itemize}
  \item Consider ord, 3.\\
$3^{1} \equiv 3(\bmod 7), 3^{2} \equiv 2(\bmod 7), 3^{3} \equiv 6(\bmod 7)$,\\
$3^{4} \equiv 4(\bmod 7), 3^{5} \equiv 5(\bmod 7), 3^{6} \equiv 1(\bmod 7)$.\\
Thus $\mathrm{ord}_{2} 3=6$. Note that Euler's Theorem guarantees that $\operatorname{ord}_{7} 3 \leqslant \phi(7)=6$.
  \item We have seen that Euler's Theovem guarantees that $\operatorname{ord}_{m} a \leqslant \phi(m)$. The next Proposition gives us the stronger statement that ordma|cl(m).
\end{itemize}

Proposition 5.1 Let $a, m \in \mathbb{Z}$ with $m>0$ and $(a, m)=1$. Then $a^{n} \equiv 1(\bmod m)$ for some positive integer $n$ if and only if ordma $\mid n$. In particular, ordma| $\phi(m)$.\\
Proof: $(\Rightarrow$ direction $)$ : Assume that $a^{n} \equiv 1(\bmod m)$ for some positive integer $n$. By the Division algorithm, there exist $q, r \in \mathbb{Z}$ such that $n=\left(\operatorname{ord} d_{m} a\right) q+r, \quad 0 \leqslant r<\operatorname{ord} d_{m} a$.

Then $a^{n}=a^{\left.\operatorname{Gor} d_{m} a\right) q+r}=\left(a^{\operatorname{ord} d_{m} a}\right)^{q} a^{r}$ $\equiv a^{r}(\bmod m)$.\\
Since $a^{n} \equiv 1(\bmod m)$, we have that $a^{r} \equiv 1(\bmod m)$.\\
Now, $O \leqslant r<o r d_{m} a$, and the definition of ord ${ }_{m}$ a as the least positive power of $a$ that is congruent to I modulom, together imply that $n=0$. Thus $n=(\operatorname{ord} m a) q$ and so ordma|n, as desired.\\
$\left(\epsilon\right.$ direction : Assume that ord ${ }_{m} a \ln$. Then there exists $b \in \mathbb{Z}$ with $b>0$ such that $n=\left(0.0 d_{m} a\right) b$.\\
Now,

$$
a^{n}=a^{(\operatorname{ord} m a) b}=\left(a^{\operatorname{ord} m a}\right)^{b} \equiv 1(\bmod m),
$$

as desired.\\
Example: By Proposition 5.) we have that ord, 2 | $\phi(7)=6$. Thus ord, 2 is $1,2,3$ or 6 . We showed previously that $\operatorname{ord}, 2=3$

\begin{itemize}
  \item By Proposition 5.1 we have that\\
$\operatorname{ord}_{13} 2 \mid \phi(13)=12$. Thus ord 132 is $1,2,3$, (4)\\
4,6 , or 12 . We have:\\
$2^{1} \equiv 2(\bmod 13), 2^{2} \equiv 4(\bmod 13), 2^{3} \equiv 8(\bmod 13)$,\\
$2^{4} \equiv 3(\bmod 13)$ and $2^{6} \equiv 12(\bmod 13)$.\\
Since we know that $2^{12} \equiv 1(\bmod 13)$ by\\
Euler's Theorom, we have that ord $132=12$\\
(without having to compute $2^{7}, 2^{8}, 2^{9}, 2^{10}$, and $\left.2^{11}(\bmod 13)\right)$.
\end{itemize}

Proposition 5.2: Let $a, m \in \mathbb{Z}$ with $m>0$ and\\
$(a, m)=1$. If $i$ and $j$ are non-negative\\
integers, then $a^{i} \equiv a^{j}(\bmod m)$ if and only\\
if $i \equiv j\left(\bmod o r d_{m} a\right)$.\\
Proof: Without loss of generality, we may assume that $i \geqslant j$.\\
$\Leftrightarrow$ direction): Assume that $a^{i} \equiv a^{d}(\bmod m)$, or, equivalently $a^{\dot{d}} a^{i-j} \equiv a^{\dot{d}}(\bmod m)$ (since $\left.i \geqslant j\right)$.

Since $(a j, m)=1$, the inverse of $a^{j}$, say (5) $\left(a^{j}\right)^{\prime}$, exists by Conollary 2.8. Multiplying both sides of the above congruence by $\left(a^{d}\right)^{\prime}$, we obtain $\left(a^{j}\right)^{\prime} a^{j} a^{i-j} \equiv\left(a^{j}\right)^{\prime} a^{j}(\bmod m)$.\\
Thus $\quad a^{i-j} \equiv 1(\bmod m)$.\\
By Proposition 5.1, we have that $\operatorname{ord}_{m} a \mid i-j$, from which $i \equiv j(\bmod \operatorname{ord} m)$, as desired.\\
$\left(\Leftarrow\right.$ direction) Assume that $i \equiv j\left(\bmod _{\text {ord }} a\right)$\\
Then $\operatorname{ord}_{m} a \mid i-j$, and so there exists $n \in \mathbb{Z}$ with $n>0$ such that $i-j=\left(\operatorname{ord} d_{m} a\right)_{n}$. Now

$$
\begin{aligned}
a^{i}=a^{\left.\operatorname{Gor} d_{m} a\right) n+j} & =a^{\left.\operatorname{Gor} d_{m} a\right) n} a^{j} \\
& =\left(a^{\operatorname{ordm} a}\right)^{n} a^{j} \\
& =a^{j}(\bmod m),
\end{aligned}
$$

as desived.\\
Example: Let $a=2$ and $m=7$. Then\\
$\operatorname{ord} 7^{2}=3$ by the previous example.

Proposition 5.2 implies that if $i$ and $j$ are non-negative integers, then $2^{i} \equiv 2^{j}(\bmod 7)$ if and only if $i \equiv j(\bmod 3)$. For example,\\
since $2000 \equiv 2(\bmod 3)$, we have that $2^{2000} \equiv 2^{2} \equiv 4(\bmod 7)$.\\
Definition: Let $r, m \geq \mathbb{Z}$ with $m>0$ and $(r, m)=1$. Then $r$ is said to be a primitive root modulo $m$ if $\operatorname{ord}_{m} n=\phi(m)$.\\
Example: 3 is a primitive root modulo 7\\
Since ord, $3=6=\phi(7)$ by a previous example.

\begin{itemize}
  \item 2 is a primitive root modulo 13 since $\operatorname{ord}_{13} 2=12=\phi(13)$ by a previous example.
  \item 2 is not a primitive root modulo 7 since $\operatorname{ord}_{2} 2=3 \neq \phi(7)$ by a previous example.\\
Example: Prove that there are no primitive\\
roots modulo 8. The reduced residues modulo
\end{itemize}

8 are 1,3,5 and 7. A straightforward check shows that $\operatorname{ord} d_{8} 1=1, \operatorname{ord} d_{8} 3=2, \operatorname{ord} d_{8} 5=2$ and $\operatorname{ord}_{8} 7=2$. Since none of the above orders are equal to $\phi(8)=4$, we have that there are no primitive roots modulo $\delta$.

\begin{itemize}
  \item Not all integers mu possess primitive roots. The Primitive Root Theorem will determire precisely which integers have primitive roots (in a later lecture).\\
Proposition 5.3: Let $r$ be a primitive root Modulo $m$. Then
\end{itemize}

$$
\left\{r, r^{2}, r^{3}, \ldots, r^{\phi(m)}\right\}
$$

is a set of reduced residues modulo $m$.\\
Remark: This result says that a primitive root (if it exists), generates the set of reduced residue classes modulo $m$.

Proof: Since $r$ is a primitive root modulo (8) $m$, we have that $(r, m)=1$, and so each element of the set is velatively prime to $m$. Since there are $\phi(m)$ such elements, it remains to show that no two elements of the set are congruent modulo m. Assume that $r^{i} \equiv r^{j}(\bmod m)$ for $1 \leq i, j \leq \phi(m)$. Then Proposition 5.2 implies that $c \equiv j(\bmod \phi(m))$, from which $i=j$ since $1 \leq 1, j \leq \phi(m)$. This Completes the proof.

Example: 3 is a primitive root modulo 7 by a previous example. One can check

$$
\begin{aligned}
& \left\{3,3^{2}, 3^{3}, 3^{4}, 3^{5}, 3^{6}\right\} \\
& =\{3,2,6,4,5,1\}(\bmod 7) .
\end{aligned}
$$

\begin{itemize}
  \item 2 is a primitive root modulo 13 by a Previous example. One can check that\\
$\left\{2,2^{2}, \cdots, 2^{12}\right\}$ is the set of reduced residues modulo 13. One can check that this set is $\{1,2,3, \ldots, 12\}(\bmod 13)$ (in some permuted order).
  \item If a primitive root modulo $m$ exists, it usually is not unique modulo $m$. At the start of next lecture we will determine the number of primitive roots there are modulo m (provided at least one exists).
\end{itemize}

Exercise: Show there are no primitive roots modulo 12.


\end{document}